\section{Introduction}\label{sec:introduction}

A minimal model program replaces a pair by successive birational models
on which its adjoint becomes more positive. For compact Kähler klt fourfold
pairs with effective rational boundary in the global Weil $\Q$-factorial
category, we prove that every maximal ordinary negative-ray program
terminates, starting on the given pair without an initial modification.
Its endpoint has nef adjoint when the initial adjoint is pseudo-effective;
otherwise it carries a projective Mori fibre space. Auxiliary generalized
pairs enter the proof, while every step of the resulting program is ordinary.

\subsection{Category and main theorem}

All varieties are reduced irreducible second-countable complex analytic
spaces. A Kähler form on a normal space has smooth strictly
plurisubharmonic local potentials in local embeddings. Projectivity of a
morphism means the existence of a relatively ample holomorphic line
bundle.

\begin{definition}\label{def:global-factoriality}\label{def:strong-factoriality}
A normal compact space $X$ is \emph{globally Weil $\Q$-factorial} if
every prime Weil divisor defined on the whole of $X$ is $\Q$-Cartier and
some positive reflexive power of its canonical sheaf $\omega_X$ is a line
bundle. It is \emph{globally strongly $\Q$-factorial} if every coherent
rank-one reflexive sheaf $\cF$ on $X$ has an invertible positive reflexive
power $\cF^{[m]}=(\cF^{\otimes m})^{**}$.
\end{definition}

The Weil convention is that of \cite[Definition~2.1]{DHY2023};
the strong convention is used in \cite[Section~2]{HX2026}.
The strong condition implies the Weil
condition. Neither definition imposes $\Q$-factoriality on every analytic
local ring. The distinction matters even for projective varieties; an
example is given in Appendix~\ref{sec:local-convention}.

We write $K_X$ additively for the canonical sheaf. When a canonical Weil
divisor is given, we keep that representative and transport it through
the program. We will also prove the intrinsic variant in which only the
canonical rational line bundle is specified. In that variant, identities
of adjoints mean isomorphisms of holomorphic line bundles after a common
positive multiple. Relative canonical divisors are defined using
compatible local canonical generators. In particular, a numerical
identity alone never specifies such a divisor.

For an effective rational boundary $\Delta$, put $D=K_X+\Delta$.
We use \emph{log discrepancies} throughout:
\[
 a(E;X,\Delta)
 =1+\operatorname{coeff}_E\bigl(K_W-p^*(K_X+\Delta)\bigr)
\]
on a smooth model $p:W\to X$ carrying $E$. A pair is klt when all these
numbers are positive. It is terminal when it is klt and they are greater
than one for every exceptional divisor over $X$. The same convention
will apply to the auxiliary real and generalized pairs.

Let $\BC(X)$ denote real Bott--Chern cohomology of closed $(1,1)$-forms
with local potentials, and let $\NA(X)$ be the closed cone of positive
currents of bidimension $(1,1)$ modulo their pairings with these classes.
A class is \emph{nef} if it belongs to the closure of the Kähler cone,
and \emph{pseudo-effective} if it contains a closed positive
$(1,1)$-current with local potentials. These terms for a rational line
bundle refer to its first Chern class. A class is \emph{big} if it
contains a Kähler current, meaning a closed current that dominates
a positive multiple of a Kähler form.
For a projective contraction
$f:X\to Z$, write $\rho(X/Z)$ for the dimension of the space of global
real line-bundle classes modulo zero degree on every contracted curve,
and put
\[
 \rho_{\mathrm{BC}}(X/Z)
 =\dim_{\R}\bigl(\BC(X)/f^*\BC(Z)\bigr).
\]

An \emph{ordinary $D$-negative step} is a projective divisorial
contraction, or the flip of a projective small contraction, of a
$D$-negative extremal ray of $\NA(X)$. The contraction has connected
fibres and a normal compact Kähler target; its relative rank is one and
$-D$ is relatively ample. For a flip the transformed adjoint is ample
over the same target. The boundary is pushed forward in a divisorial
step and strictly transformed in a flip. A \emph{Mori fibre space} has
the same contraction properties and a base of strictly smaller
dimension.

\begin{theorem}[Termination of ordinary programs]\label{thm:finite}
Let $X$ be a normal globally Weil $\Q$-factorial compact Kähler
fourfold, with a canonical Weil divisor $K_X$, and let $\Delta\geq0$
be a rational divisor such that $(X,\Delta)$ is klt. Then every maximal
ordinary negative-ray program starting on $(X,\Delta)$ terminates.
Its birational steps form a finite sequence
\[
 (X,\Delta)=(X_0,\Delta_0)\dashrightarrow\cdots
 \dashrightarrow(X_n,\Delta_n).
\]
Every $X_i$ is normal, compact Kähler and globally Weil
$\Q$-factorial; every $(X_i,\Delta_i)$ is klt, and
$D_i=K_{X_i}+\Delta_i$ is $\Q$-Cartier. The canonical divisors are
the transforms of the specified datum. The following alternatives hold.
\begin{enumerate}[label=\textup{(\roman*)}]
\item If $D=K_X+\Delta$ is pseudo-effective, then $D_n$ is nef.
The composite $\phi:X\dashrightarrow X_n$ extracts no prime divisor,
and
\[
 a(E;X,\Delta)\leq a(E;X_n,\Delta_n)
\]
for every prime divisor over the two models, with strict inequality
for every prime on $X$ contracted by $\phi$.
\item If $D$ is not pseudo-effective, there is a projective surjective
morphism $g:X_n\to S$ with connected fibres, where $S$ is normal
compact Kähler, $\dim S<4$, $\rho(X_n/S)=1$, and $-D_n$ is
$g$-ample.
\end{enumerate}
For each negative contraction, including $g$, the relative
Bott--Chern dimension is one. If $X$ is globally strongly
$\Q$-factorial, so is every $X_i$.

The same conclusions hold in the intrinsic formulation with the
canonical rational line bundle in place of a chosen canonical Weil
divisor. In either formulation the first model is exactly $X$.
\end{theorem}

The endpoint conclusion in the pseudo-effective case is nefness; no
abundance assertion is made.

\subsection{Prior work and proof ingredients}

Höring--Peternell established minimal models for non-uniruled compact
Kähler threefolds~\cite{HP2016}. For fourfolds,
Das--Hacon--Păun proved the compact klt result when the adjoint is
$\Q$-linearly equivalent to an effective divisor~\cite{DHP2024}.
Fujino's analytic relative MMP supplies ordinary cone, base-point-free,
and flip theorems for projective morphisms~\cite{Fujino2022}.
The generalized cone theorem of Hacon--Xie supplies the analytic
negative rays and their finite form when the boundary plus nef part is
big~\cite[Theorem~1.3]{HX2026}.

For the supporting contractions needed here, Section~\ref{sec:contractions}
proves a bounded-dimensional base-point-free statement. Its induction
uses ordinary rational programs in dimensions at most three, all with
global line-bundle polarizations. A divisorial negative-part comparison
then supplies descent from a selected lower-dimensional good model.
The ordinary rational adjoint detects each fourfold negative ray and
gives a global algebra for its flip.

The companion theorem on generalized log canonical flips
\cite[Theorem~1.1]{OAI:Termination-of-generalized-log-canonical-flips-on-compact-Kahler-fourfolds-October-7-2026}
proves termination for the resulting existing rational flip sequences.
Compact analytic cycle classes bound the number of divisorial steps.
The companion uses the independent low-place, relative-program,
extraction, and special-termination results proved below; the supporting
contraction argument does not use the finite-program theorem itself.

Two ingredients have uses beyond this deduction. First, actual line
bundles that are numerically trivial over the relevant birational
contractions descend without taking an additional power. Second,
ordinary terminal fourfold flips terminate with effective \emph{real}
boundary in our compact Kähler setting. The proof adapts the
discrepancy and homology methods of Kawamata--Matsuda--Matsuki and
Fujino, including the treatment of low-discrepancy places in his addendum
\cite{KMM1987,Fujino2004,Fujino2005Addendum}.

\subsection{Proof outline}

Sections~\ref{sec:analytic}--\ref{sec:counts} establish the common
analytic, birational and termination tools. The proof of
Theorem~\ref{thm:finite} then has three stages.
\begin{enumerate}[label=\textup{\arabic*.}]
\item The cone theorem gives a supporting nef class for each ordinary
negative ray. Once its contraction is constructed, the ordinary
rational adjoint provides one global relative polarization and one
global flip algebra.
\item A dimension induction proves the supporting contraction
statement through dimension four. Its lower-dimensional ordinary
programs are projective from their first step, and a negative-part
comparison gives descent without lifting a generalized program from
the lower-dimensional base.
\item The companion flip-termination theorem excludes an infinite
flip tail, while compact cycle counts bound the divisorial steps.
Section~\ref{sec:endpoints} identifies the endpoint from
pseudo-effectivity and gives the required discrepancy inequalities.
\end{enumerate}
All auxiliary constructions are separate from the ordinary sequence,
which begins on $X$ itself.

The rational relative constructions in Section~\ref{sec:relative}
feed threefold termination and the dimension induction for special
termination. These results support both the contraction induction and
the companion flip-termination theorem. Their nef data are actual rational
line bundles on higher models, denoted by $\mathbf M$, distinct from the
possibly transcendental nef classes used in the contraction argument.

\begin{figure}[htbp]
\centering
\begin{tikzpicture}[
 box/.style={draw=linkblue!65,rounded corners=2pt,align=center,
 fill=linkblue!3,text width=5.4cm,inner sep=7pt,font=\small},
 shared/.style={box,text width=11.8cm},
 edge/.style={-{Stealth[length=2mm]},draw=linkblue!80,thick}]
\node[shared] (common) {Analytic contractions and exact bundle descent\\
Compact cycle counts and ordinary terminal fourfold termination};
\node[box,below=9mm of common.south west,anchor=north west] (relative)
 {Rational relative programs\\
Threefold and special termination};
\node[box,below=9mm of common.south east,anchor=north east] (supporting)
 {Supporting contractions\\
Bounded-dimensional induction};
\node[box,below=8mm of relative] (companion)
 {Companion flip termination\\
Compact cycle counts};
\node[box,below=8mm of supporting] (ordinary)
 {Global line-bundle detectors\\
Ordinary projective steps};
\node[shared,below=8mm of companion.south west,anchor=north west] (finite)
 {Every maximal ordinary negative-ray program terminates\\
Theorem~\ref{thm:finite}};
\draw[edge] (common.south -| relative.north)--(relative.north);
\draw[edge] (common.south -| supporting.north)--(supporting.north);
\draw[edge] (relative)--(supporting);
\draw[edge] (relative)--(companion);
\draw[edge] (supporting)--(ordinary);
\draw[edge] (companion.south)--(finite.north -| companion.south);
\draw[edge] (ordinary.south)--(finite.north -| ordinary.south);
\end{tikzpicture}
\caption{Dependencies in the proof of Theorem~\ref{thm:finite}. The
relative and lower-dimensional results support the contraction induction
and the companion flip-termination theorem.}
\label{fig:dependencies}
\end{figure}

